\clearpage\selectlanguage{german}\fancyhead[L]{\small ZUR THEORIE DER NEUEN TRANSSCENDENTEN. III.}\pdfbookmark[0]{Zur Theorie der neuen Transscendenten III}{transscendentes-iii}\SourceLeaf{455}
\begin{center}
[III.]\par
\textbf{[ZUR THEORIE DER NEUEN TRANSSCENDENTEN.]}\par
\rule{13mm}{0.3pt}
\end{center}
Die Theoreme in Beziehung auf diejenigen Reihen und unendlichen Producte, welche zu der Theorie der Arithmetisch Geometrischen Mittel gehören,
ordnen wir so:
\[
\begin{array}{r@{\qquad}l}
1.&1-x\mathbin{.}1-xx\mathbin{.}1-x^3\mathbin{.}1-x^4\mathinner{.\,.\,.}=[x]\\[7pt]
2.&1+x\mathbin{.}1+xx\mathbin{.}1+x^3\mathbin{.}1+x^4\mathinner{.\,.\,.}=\dfrac{[xx]}{[x]}\\[11pt]
3.&1-x\mathbin{.}1-x^3\mathbin{.}1-x^5\mathbin{.}1-x^7\mathinner{.\,.\,.}=\dfrac{[x]}{[xx]}\\[11pt]
4.&1+x\mathbin{.}1+x^3\mathbin{.}1+x^5\mathbin{.}1+x^7\mathinner{.\,.\,.}=\dfrac{[xx]^2}{[x][x^4]}\\[7pt]
5.&[-x]=\dfrac{[xx]^3}{[x][x^4]}
\end{array}
\]
\[
\begin{gathered}6.\quad1+xy\mathbin{.}1+x^3y\mathbin{.}1+x^5y\mathinner{.\,.\,.}
\\[8pt]\quad{}1+xy^{-1}\mathbin{.}1+x^3y^{-1}\mathbin{.}1+x^5y^{-1}\mathinner{.\,.\,.}\end{gathered}
\]
evolvitur in seriem
\[
\begin{gathered}Fx\mathbin{.}\{1+(y+y^{-1})x+(y^2+y^{-2})x^4\\[8pt]\quad{}+(y^3+y^{-3})x^9+\mathinner{.\,.\,.}\}\end{gathered}
\]
\[
\begin{gathered}7.\quad\text{also}\quad
Fx=\frac{(1-x)^2(1-x^3)^2(1-x^5)^2\mathinner{.\,.\,.}}
{1-2x+2x^4-2x^9+\mathinner{.\,.\,.}}
\\[8pt]\quad{}=\frac{[x]^2}{[xx]^2}\mathbin{.}
\frac{1}{1-2x+2x^4-2x^9+\mathinner{.\,.\,.}}\end{gathered}
\]
\SourceLeaf{456}
\[
\begin{gathered}8.\quad Fx=
\frac{1+xx\mathbin{.}1+x^6\mathbin{.}1+x^{10}\mathinner{.\,.\,.}}
{1-2x^4+2x^{16}-2x^{36}+\mathinner{.\,.\,.}}
\\[8pt]\quad{}=\frac{[x^4]^2}{[xx][x^8]}\mathbin{.}
\frac{1}{1-2x^4+2x^{16}-\mathinner{.\,.\,.}}\end{gathered}
\]
\[
\begin{gathered}9.\quad [xx]Fx=[x^8]Fx^4=[x^{32}]Fx^{16}\\[8pt]\quad{}=[x^{128}]Fx^{64}=\mathrm{etc.}=1\end{gathered}
\]
\[
\begin{gathered}10.\quad1-2x+2x^4-\mathinner{.\,.}
=\frac{[x]^2}{[xx]}
\\[8pt]\quad{}=\frac{1-x\mathbin{.}1-xx\mathbin{.}1-x^3\mathinner{.\,.}}
{1+x\mathbin{.}1+xx\mathbin{.}1+x^3\mathinner{.\,.}}\end{gathered}
\]
\[
\begin{gathered}11.\quad1+2x+2x^4+\mathinner{.\,.}
=\frac{[xx]^5}{[x]^2[x^4]^2}
\\[8pt]\quad{}=\frac{1+x\mathbin{.}1-xx\mathbin{.}1+x^3\mathbin{.}1-x^4\mathinner{.\,.}}
{1-x\mathbin{.}1+xx\mathbin{.}1-x^3\mathbin{.}1+x^4\mathinner{.\,.}}\end{gathered}
\]
Andere Beweise dieser Sätze.

Wenn man in 6 statt $y$, $xy$ schreibt, so wird
\[
\begin{gathered}
1+\frac{1}{yy}\mathbin{.}1+xxyy\mathbin{.}1+x^4yy\mathbin{.}1+x^6yy\mathinner{.\,.\,.}
\\[8pt]\quad{}1+xxy^{-2}\mathbin{.}1+x^4y^{-2}\mathbin{.}1+x^6y^{-2}\mathinner{.\,.}\\[8pt]
\qquad=\frac{1}{[xx]}\{(1+y^{-2})+(y^2+y^{-4})xx+(y^4+y^{-6})x^6+\mathinner{.\,.}\}
\end{gathered}
\]
oder
\[
\begin{gathered}
12.\quad{}y+\frac{1}{y}\mathbin{.}1+xxyy\mathbin{.}1+x^4yy\mathbin{.}1+x^6yy\mathinner{.\,.\,.}
\\[8pt]\quad{}1+xxy^{-2}\mathbin{.}1+x^4y^{-2}\mathbin{.}1+x^6y^{-2}\mathinner{.\,.}\\[8pt]
\qquad=\frac{1}{[xx]x^{\frac14}}\{(y+y^{-1})x^{\frac14}+(y^3+y^{-3})x^{\frac94}+\mathinner{.\,.}\}
\end{gathered}
\]
Anderer Beweis
\[
13.\quad x^{\frac14}\frac{[x^4]^2}{[xx]}=x^{\frac14}+x^{\frac94}+\mathinner{.\,.\,.}
\]
oder
\[
\begin{gathered}14.\quad\frac{[xx]^2}{[x]}=1+x+x^3+x^6+x^{10}+\mathinner{.\,.}
\\[8pt]\quad{}=\frac{1-xx\mathbin{.}1-x^4\mathinner{.\,.}}{1-x\mathbin{.}1-x^3\mathinner{.\,.}}\end{gathered}
\]
Anderer Beweis.
\[
\begin{gathered}15.\quad(1-2x+2x^4-\mathinner{.\,.})\quad(1+2x+2x^4+\mathinner{.\,.})
\\[8pt]=(1-2xx+2x^8-\mathinner{.\,.})^2\end{gathered}
\]
\[
\begin{gathered}16.\quad(1-2x+2x^4-\mathinner{.\,.})^2+(1+2x+2x^4+\mathinner{.\,.})
\\[8pt]=2(1+2xx+2x^8+\mathinner{.\,.})^2\end{gathered}
\]
\[
\begin{gathered}17.\quad(1+2x+2x^4+\mathinner{.\,.})^2(x^{\frac14}+x^{\frac94}+\mathinner{.\,.})
\\[8pt]=(x^{\frac18}+x^{\frac98}+\mathinner{.\,.})^2\end{gathered}
\]
\[
\begin{gathered}18.\quad(1+2x+2x^4+\mathinner{.\,.})^2+(2x^{\frac14}+2x^{\frac94}+\mathinner{.\,.})^2
\\[8pt]=(1+2x^{\frac12}+2x^2+\mathinner{.\,.})^2\end{gathered}
\]
\[
\begin{gathered}19.\quad(1+2x+2x^4+\mathinner{.\,.})^4
\\[8pt]=(1-2x+2x^4-\mathinner{.\,.})^4+(2x^{\frac14}+2x^{\frac14}+\mathinner{.\,.})^4\end{gathered}
\]
\SourceLeaf{457}
Anwendung auf arithm.\ geom.\ Mittel.
\[
\begin{gathered}
20.\quad a=h(1+2x+2x^4+\mathinner{.\,.})^2\\[6pt]b=h(1-2x+2x^4-\mathinner{.\,.})^2\\[7pt]
a'=h(1+2xx+2x^8+\mathinner{.\,.})^2\\[6pt]b'=h(1-2xx+2x^8-\mathinner{.\,.})^2\\[7pt]
a''=h(1+2x^4+2x^{16}+\mathinner{.\,.})^2\\[6pt]b''=h(1-2x^4+2x^{16}-\mathinner{.\,.})^2\\[7pt]
a'''=h(1+2x^8+2x^{32}+\mathinner{.\,.})^2\\[6pt]b'''=h(1-2x^8+2x^{32}-\mathinner{.\,.})^2\\[5pt]
\qquad\mathrm{etc.}\qquad\mathrm{etc.}
\end{gathered}
\]
\[
\begin{gathered}
c=\surd(aa-bb)\\[8pt]
c'=\surd(a'a'-b'b')\\[7pt]=\tfrac12(a-b)=\dfrac{cc}{4a'}=\dfrac{cc}{4a'},
\\[7pt]\surd\dfrac{c'}{4h}=\dfrac{c}{4}\dfrac{1}{\surd a'h}\\[12pt]
c''=\surd(a''a''-b''b'')\\[7pt]=\tfrac12(a'-b')=\dfrac{c'c'}{4a''}=\dfrac{c^4}{64a'a'a''},
\\[7pt]{}^{4}\!\surd\dfrac{c''}{4h}=\dfrac{c}{4{a'}^{\frac12}{a''}^{\frac14}h^{\frac14}}\\[12pt]
c'''=\surd(a'''a'''-b'''b''')\\[7pt]=\tfrac12(a''-b'')=\dfrac{c''c''}{4a'''}=\dfrac{c^8}{2^{14}{a'}^4{a''}^2a'''},
\\[7pt]{}^{8}\!\surd\dfrac{c'''}{4h}=\dfrac{c}{4{a'}^{\frac12}{a''}^{\frac14}{a'''}^{\frac18}h^{\frac18}}\\[7pt]
\qquad\mathrm{etc.}\\[7pt]\mathrm{etc.}
\end{gathered}
\]
\[
\begin{gathered}21.\quad x^{\frac12}=\frac{c}{4{a'}^{\frac12}{a''}^{\frac14}{a'''}^{\frac18}\mathinner{.\,.}}
\\[8pt]\quad{}=\frac{c}{4a'}\Bigl[\frac{a'}{a''}\Bigr]^{\frac12}
\Bigl[\frac{a''}{a'''}\Bigr]^{\frac14}
\Bigl[\frac{a'''}{a''''}\Bigr]^{\frac18}\ \mathrm{etc.}\end{gathered}
\]
\[
22.\quad x=\frac{a-b}{8a''}\left[\frac{a''}{a'''}\right]^{\frac12}
\left[\frac{a'''}{a''''}\right]^{\frac14}\ \mathrm{etc.}
\]
Setzt man in 6 statt $x$, $x^3$ und statt $y$, $-x$, so wird
\[
\begin{gathered}
1-x^4\mathbin{.}1-x^{10}\mathbin{.}1-x^{16}\mathinner{.\,.\,.}
\\[8pt]\quad{}1-x^2\mathbin{.}1-x^8\mathbin{.}1-x^{14}\mathinner{.\,.\,.}\\[8pt]
\qquad=\frac{1}{[x^6]}\{1-xx-x^4+x^{10}+x^{14}-x^{24}-x^{30}+\mathinner{.\,.}\}
\end{gathered}
\]
oder
\[
23.\quad[x]=1-x-xx+x^5+x^7-x^{12}-x^{15}+\mathrm{etc.}
\]
Anderer Beweis.

Man setze in 6 statt $x$, $x^3$ und statt $y$, $+x$, so wird

\SourceLeaf{458}
\[
\begin{gathered}
24.\quad 1+x+xx+x^{5}+x^{7}+x^{12}+x^{15}+\mathrm{etc.}
 \\[8pt]\quad{}= [x^{3}]\mathbin{.}1+x\mathbin{.}1+xx\mathbin{.}1+x^{4}\mathbin{.}1+x^{5}\mathinner{.\,.}\\[8pt]
 \\[8pt]\quad{}=\frac{[xx][x^{3}]^{2}}{[x][x^{6}]}
\end{gathered}
\]
Man setze in 6 statt $x$, $x^{3}$ und statt $y$, $xxy$, so wird
\[
\begin{gathered}
1+\frac{x}{y}+x^{5}y+\frac{x^{8}}{yy}+x^{16}yy+\mathrm{etc.}\\[8pt]
\qquad=1+\frac{x}{y}\mathbin{.}1+\frac{x^{7}}{y}\mathbin{.}1+\frac{x^{13}}{y}\mathinner{.\,.\,.}
  \\[8pt]\quad{}1+x^{5}y\mathbin{.}1+x^{11}y\mathbin{.}1+x^{17}y\mathinner{.\,.\,.}[x^{6}]
\end{gathered}
\]
oder statt $x$, $x^{3}$ gesetzt
\[
\begin{gathered}
x+\frac{x^{4}}{y}+x^{16}y+\frac{x^{25}}{yy}+x^{49}yy+\mathrm{etc.}\\[8pt]
\qquad=x\mathbin{.}1+\frac{x^{3}}{y}\mathbin{.}1+\frac{x^{21}}{y}\mathbin{.}1+\frac{x^{39}}{y}\mathinner{.\,.\,.}
  \\[8pt]\quad{}1+x^{15}y\mathbin{.}1+x^{33}y\mathbin{.}1+x^{51}y\mathinner{.\,.\,.}[x^{18}]
\end{gathered}
\]
also
\[
\begin{gathered}
25.\quad x+x^{4}+x^{16}+x^{25}+x^{49}+\mathinner{.\,.}\\[8pt]
 =x\mathbin{.}1+x^{3}\mathbin{.}1+x^{15}\mathbin{.}1+x^{21}\\[8pt]\quad{}\mathbin{.}1+x^{33}\mathbin{.}1+x^{39}\mathinner{.\,.\,.}[x^{18}]
 \\[8pt]\quad{}=x\frac{[x^{6}]^{2}[x^{9}][x^{36}]}{[x^{3}][x^{12}][x^{18}]}
\end{gathered}
\]
\[
\begin{gathered}
26.\quad x-x^{4}-x^{16}+x^{25}+x^{49}-\mathrm{etc.}
 \\[8pt]\quad{}=x\mathbin{.}1-x^{3}\mathbin{.}1-x^{15}\mathbin{.}1-x^{21}\mathbin{.}1-x^{33}\mathinner{.\,.}[x^{18}]\\[8pt]
 \\[8pt]\quad{}=x\frac{[x^{3}][x^{18}]^{2}}{[x^{6}][x^{9}]}
\end{gathered}
\]
Da nun
\[
\begin{gathered}\frac{x\mathbin{.}1-x^{3}\mathbin{.}1-x^{15}\mathbin{.}1-x^{21}\mathbin{.}1-x^{33}\mathinner{.\,.\,.}[x^{18}]}
     {x\mathbin{.}1-x^{3}\mathbin{.}1-x^{9}\mathbin{.}1-x^{15}\mathbin{.}1-x^{21}\mathinner{.\,.\,.}}
 \\[8pt]\quad{}=\frac{[x^{18}]^{2}}{[x^{9}]}
\\[8pt]\quad{} =1+x^{9}+x^{27}+x^{54}+\mathinner{.\,.}\end{gathered}
\]
so ist
\[
\begin{gathered}
27.\quad x-x^{4}-x^{16}+x^{25}+x^{29}-\mathinner{.\,.}\\[8pt]
 =1-x^{3}\mathbin{.}1-x^{9}\mathbin{.}1-x^{15}\mathinner{.\,.\,.}\frac{1}{x^{\frac18}}
 \\[8pt]\quad{}\{x^{\frac98}+x^{\frac{81}{8}}+x^{\frac{225}{8}}+x^{\frac{441}{8}}+\mathrm{etc.}\}
\end{gathered}
\]
ferner folgt aus 24, wenn man statt $x$, $x^{3}$ setzt, weil
\[
\begin{gathered}\frac{1-2x^{9}+2x^{36}-2x^{81}+\mathinner{.\,.\,.}}
     {1-x^{3}\mathbin{.}1-x^{9}\mathbin{.}1-x^{15}\mathinner{.\,.\,.}}
 \\[8pt]\quad{}=\frac{[x^{6}][x^{9}]^{2}}{[x^{3}][x^{18}]}
\\[8pt]\quad{} =1+x^{3}+x^{6}+x^{15}+x^{21}+\mathinner{.\,.}\end{gathered}
\]
\[
\begin{gathered}28.\quad 1-2x^{9}+2x^{36}-2x^{81}+\mathinner{.\,.}
\\[8pt]\quad{} =1-x^{3}\mathbin{.}1-x^{9}\mathbin{.}1-x^{15}\mathinner{.\,.\,.}\frac{1}{x^{\frac18}}
 \\[8pt]\quad{}\{x^{\frac18}+x^{\frac{25}{8}}+x^{\frac{49}{8}}+\mathrm{etc.}\}\end{gathered}
\]
also aus der Verbindung von 27 und 28
\SourceLeaf{459}
\[
\begin{gathered}
29.\quad 1-2x+2x^{4}-2x^{9}+2x^{16}-\mathinner{.\,.}\\[8pt]
 =1-x^{3}\mathbin{.}1-x^{9}\mathbin{.}1-x^{15}\mathinner{.\,.\,.}\frac{1}{x^{\frac18}}
 \\[8pt]\quad{}\{x^{\frac18}-2x^{\frac98}+x^{\frac{25}{8}}+x^{\frac{49}{8}}-2x^{\frac{81}{8}}-\mathrm{etc.}\}
\end{gathered}
\]
ferner folgt aus 23
\[
x^{\frac18}-x^{\frac{25}{8}}-x^{\frac{49}{8}}+x^{\frac{121}{8}}+\mathinner{.\,.}
 =x^{\frac18}[x^{3}]
\]
also
\[
\begin{gathered}
30.\quad 1-2x^{3}+2x^{12}-2x^{27}+2x^{48}-\mathinner{.\,.}\\[8pt]
 =1-x^{3}\mathbin{.}1-x^{9}\mathbin{.}1-x^{15}\mathinner{.\,.\,.}\frac{1}{x^{\frac18}}
 \\[8pt]\quad{}\{x^{\frac18}-x^{\frac{25}{8}}-x^{\frac{49}{8}}+x^{\frac{121}{8}}+x^{\frac{169}{8}}-\mathinner{.\,.}\}
\end{gathered}
\]
Aus der Summation von 29 und 30
\[
\begin{gathered}
31.\quad \{1-2x+2x^{4}-2x^{9}+\mathinner{.\,.}\}\\[8pt]\quad{}+\{1-2x^{3}+2x^{12}-2x^{27}+\mathinner{.\,.}\}\\[8pt]
 =1-x^{3}\mathbin{.}1-x^{9}\mathbin{.}1-x^{15}\mathinner{.\,.\,.}\frac{2}{x^{\frac18}}
 \\[8pt]\quad{}\{x^{\frac18}-x^{\frac98}-x^{\frac{81}{8}}+x^{\frac{121}{8}}+x^{\frac{169}{8}}-\mathinner{.\,.}\}
\end{gathered}
\]
Aus der Summation von 28 und 30
\[
\begin{gathered}
32.\quad \{1-2x^{3}+2x^{12}-2x^{27}+\mathinner{.\,.}\}\\[8pt]\quad{}+\{1-2x^{9}+2x^{36}-2x^{81}+\mathinner{.\,.}\}\\[8pt]
 =1-x^{3}\mathbin{.}1-x^{9}\mathbin{.}1-x^{15}\mathinner{.\,.\,.}\frac{2}{x^{\frac18}}
 \\[8pt]\quad{}\{x^{\frac18}+x^{\frac{121}{8}}+x^{\frac{169}{8}}+x^{\frac{529}{8}}+\mathinner{.\,.}\}
\end{gathered}
\]
Setzt man in 6 statt $x$, $x^{6}$ und statt $y$, $xy$, so wird
\[
\begin{gathered}
1+x^{7}y\mathbin{.}1+x^{19}y\mathbin{.}1+x^{31}y\mathinner{.\,.\,.}\\[8pt]\quad{}1+\frac{x^{5}}{y}\mathbin{.}1+\frac{x^{17}}{y}\mathinner{.\,.\,.}[x^{12}]\\[8pt]
\qquad=\\[8pt]\quad{}1+\frac{x^{5}}{y}+x^{7}y+\frac{x^{22}}{yy}+x^{26}yy+\mathinner{.\,.}
\end{gathered}
\]
Man hat demnach die Zerlegungen in Factoren
\[
\begin{gathered}
33.\quad (1-2x^{3}+2x^{12}-2x^{27}+\mathinner{.\,.})\\[8pt]\quad{}+(1-2x^{9}+2x^{36}-2x^{81}+\mathinner{.\,.})\\[8pt]
 =2[x^{36}]\mathbin{.}1-x^{3}\mathbin{.}1-x^{9}\mathbin{.}1-x^{15}\mathinner{.\,.\,.}\\[8pt]\quad{}
 1+x^{15}\mathbin{.}1+x^{21}\mathbin{.}1+x^{51}\mathbin{.}1+x^{57}\mathinner{.\,.\,.}
\end{gathered}
\]
\[
\begin{gathered}
34.\quad (1-2x+2x^{4}-2x^{9}+\mathinner{.\,.})\\[8pt]\quad{}+(1-2x^{3}+2x^{12}-2x^{27}+\mathinner{.\,.})\\[8pt]
 =2[x^{12}]\mathbin{.}1-x\mathbin{.}1-x^{3}\mathbin{.}1-x^{5}\mathbin{.}1-x^{7}\mathinner{.\,.\,.}\\[8pt]\quad{}
 1+x^{5}\mathbin{.}1+x^{7}\mathbin{.}1+x^{17}\mathbin{.}1+x^{19}\mathinner{.\,.\,.}
\end{gathered}
\]
\[
\begin{gathered}
35.\quad (1+2x^{3}+2x^{12}+2x^{27}+\mathinner{.\,.})\\[8pt]\quad{}+(1+2x^{9}+2x^{36}+2x^{81}+\mathinner{.\,.})\\[8pt]
 =2[x^{36}]\mathbin{.}1+x^{3}\mathbin{.}1+x^{9}\mathbin{.}1+x^{15}\mathinner{.\,.\,.}\\[8pt]\quad{}
 1-x^{15}\mathbin{.}1-x^{21}\mathbin{.}1-x^{51}\mathinner{.\,.\,.}
\end{gathered}
\]
\[
\begin{gathered}
36.\quad (1+2x+2x^{4}+2x^{9}+\mathinner{.\,.})\\[8pt]\quad{}+(1+2x^{3}+2x^{12}+2x^{27}+\mathinner{.\,.})\\[8pt]
 =2[x^{12}]\mathbin{.}1+x\mathbin{.}1+x^{3}\mathbin{.}1+x^{5}\mathbin{.}1+x^{7}\mathinner{.\,.\,.}\\[8pt]\quad{}
 1-x^{5}\mathbin{.}1-x^{7}\mathbin{.}1-x^{17}\mathbin{.}1-x^{19}\mathinner{.\,.\,.}
\end{gathered}
\]

\SourceLeaf{460}

Hieraus ergeben sich zugleich die Factoren des letzten Theils in 31.

Aus der Subtraction von 28 und 30
\[
\begin{gathered}
37.\quad (1-2x^9+2x^{36}-2x^{81}+ .\,.)\\[8pt]\quad{}-(1-2x^3+2x^{12}- .\,.)\\[8pt]
\qquad=2\,.\,1-x^3\,.\,1-x^9\,.\,1-x^{15}\,.\,.\,
\\[8pt]\quad{}\frac{1}{x^{\frac18}}\bigl(x^{\frac{25}{8}}+x^{\frac{49}{8}}+x^{\frac{289}{8}}+x^{\frac{361}{8}}+ .\,.\bigr)
\end{gathered}
\]
Setzt man in 6 statt $x$, $x^6$ und statt $y$, $x^5y$, so wird
\[
\begin{gathered}
1+x^{11}y\,.\,1+x^{23}y\,.\,1+x^{35}y\,.\,.\,.\,
\\[8pt]\quad{}1+\frac{x}{y}\,.\,1+\frac{x^{13}}{y}\,.\,1+\frac{x^{25}}{y}\,.\,.\,.\,.\,[x^{12}]\\[8pt]
\qquad=\\[8pt]\quad{}1+\frac{x}{y}+x^{11}y+\frac{x^{14}}{yy}+x^{34}yy+ .\,.\,.
\end{gathered}
\]
Also die Zerlegung in Factoren
\[
\begin{gathered}
38.\quad (1-2x^5+2x^{36}-2x^{81}+ .\,.)\\[8pt]\quad{}-(1-2x^3+2x^{12}- .\,.)\\[8pt]
=2x^3[x^{36}]\,.\,1+x^3\,.\,1+x^{33}\,.\,1+x^{39}\,.\,1+x^{69}\,.\,.\,.\,\\[8pt]\quad{}
1-x^3\,.\,1-x^9\,.\,1-x^{15}\,.\,.\,.
\end{gathered}
\]
\[
\begin{gathered}
39.\quad (1-2x^3+2x^{12}-2x^{27}+ .\,.)\\[8pt]\quad{}-(1-2x+2x^4- .\,.)\\[8pt]
=2x[x^{12}]\,.\,1+x\,.\,1+x^{11}\,.\,1+x^{13}\,.\,1+x^{23}\,.\,.\,.\,\\[8pt]\quad{}
1-x\,.\,1-x^3\,.\,1-x^5\,.\,.\,.
\end{gathered}
\]
\[
\begin{gathered}
40.\quad (1+2x^3+2x^{12}+2x^{27}+ .\,.)\\[8pt]\quad{}-(1+2x^9+2x^{36}+ .\,.)\\[8pt]
=2x^3[x^{36}]\,.\,1-x^3\,.\,1-x^{33}\,.\,1-x^{39}\,.\,1-x^{69}\,.\,.\,.\,\\[8pt]\quad{}
1+x^3\,.\,1+x^9\,.\,1+x^{15}\,.\,.\,.
\end{gathered}
\]
\[
\begin{gathered}
41.\quad (1+2x+2x^4+2x^9+ .\,.)\\[8pt]\quad{}-(1+2x^3+2x^{12}+ .\,.)\\[8pt]
=2x[x^{12}]\,.\,1-x\,.\,1-x^{11}\,.\,1-x^{13}\,.\,1-x^{23}\,.\,.\,.\,\\[8pt]\quad{}
1+x\,.\,1+x^3\,.\,1+x^5\,.\,.\,.
\end{gathered}
\]
Aus der Subtraction von 29 und 30 folgt
\[
\begin{gathered}
42.\quad (1-2x^3+2x^{12}-2x^{27}+ .\,.)\\[8pt]\quad{}-(1-2x+2x^4- .\,.)\\[8pt]
\qquad=2\,.\,1-x^3\,.\,1-x^9\,.\,1-x^{15}\,.\,.\,.\,
\\[8pt]\quad{}\frac{1}{x^{\frac18}}\bigl(x^{\frac98}-x^{\frac{25}{8}}-x^{\frac{49}{8}}+x^{\frac{81}{8}}+x^{\frac{225}{8}}- .\,.\bigr)
\end{gathered}
\]
Woraus die Zerlegung des letzten Gliedes dieser Gleichung in Factoren folgt.

Aus der Multiplication von 34 und 39 folgt
\[
\begin{gathered}
43.\quad (1-2x^3+2x^{12}-2x^{27}+ .\,.)^2\\[8pt]\quad{}-(1-2x+2x^4- .\,.)^2\\[8pt]
\quad=4x[x^{12}]^2\,.\,(1-x)^2(1-x^3)^2(1-x^5)^2\,.\,.\,.\,\\[8pt]\quad{}
1+x\,.\,1+x^5\,.\,1+x^7\,.\,1+x^{11}\,.\,.\,.\\[8pt]
\quad=4x\frac{[x^{12}]^2[x]^2}{[xx]^2}\,.\,
\frac{[xx]^2[x^3][x^{12}]}{[x][x^4][x^6]^2}
\\[8pt]\quad{}=4x\frac{[x][x^3][x^{12}]^3}{[x^4][x^6]^2}
\end{gathered}
\]
Ebenso aus der Multiplication von 36 und 41

\SourceLeaf{461}
\[
\begin{gathered}
44.\quad (1+2x+2x^4+ .\,.)^2\\[8pt]\quad{}-(1+2x^3+2x^{12}+2x^{27}+ .\,.)^2\\[8pt]
\quad=4x[x^{12}]^2\,.\,(1+x)^2(1+x^3)^2(1+x^5)^2\,.\,.\,.\,\\[8pt]\quad{}
1-x\,.\,1-x^5\,.\,1-x^7\,.\,1-x^{11}\,.\,.\,.\\[8pt]
\quad=4x\frac{[x^{12}]^2[xx]^4}{[x]^2[x^4]^2}\,.\,
\frac{[x][x^6]}{[xx][x^3]}
\\[8pt]\quad{}=4x\frac{[xx]^3[x^6][x^{12}]^2}{[x][x^3][x^4]^2}
\end{gathered}
\]
Also der Quotient
\[
\begin{gathered}
45.\quad \frac{(1-2x^3+2x^{12}-2x^{27}+ .\,.)^2-(1-2x+2x^4- .\,.)^2}
{(1+2x^3+2x^{12}+2x^{27}+ .\,.)^2-(1+2x+2x^4+ .\,.)^2}\\[8pt]
\qquad=-\Bigl(\frac{1-x}{1+x}\Bigr)^3
\Bigl(\frac{1-x^3}{1+x^3}\Bigr)
\Bigl(\frac{1-x^5}{1+x^5}\Bigr)^3
\\[8pt]\quad{}\Bigl(\frac{1-x^7}{1+x^7}\Bigr)^3
\Bigl(\frac{1-x^9}{1+x^9}\Bigr)\,.\,.\,.\\[8pt]
\qquad=\frac{[x]^2[x^3]^2[x^4][x^{12}]}{[xx]^3[x^6]^3}
\end{gathered}
\]
und das Product
\[
\begin{gathered}
45^{\mathrm b}.\quad \bigl\{(1-2x^3+2x^{12}- .\,.)^2-(1-2x+2x^4- .\,.)^2\bigr\}\\[8pt]
\quad\times\bigl\{(1+2x^3+2x^{12}+ .\,.)^2-(1+2x+2x^4+ .\,.)^2\bigr\}\\[8pt]
\qquad=-16xx\frac{[xx]^3[x^{12}]^5}{[x^4]^3[x^6]}
\end{gathered}
\]
Aus $28+i\,30$ folgt
\[
\begin{gathered}
46.\quad (1-2x^9+2x^{36}- .\,.)\\[8pt]\quad{}-i(1-2x^3+2x^{12}- .\,.)\\[8pt]
\qquad=1-x^3\,.\,1-x^9\,.\,1-x^{15}\,.\,.\,.\,
\\[8pt]\quad{}\frac{1-i}{x^{\frac18}}\bigl(x^{\frac18}+ix^{\frac{25}{8}}+ix^{\frac{49}{8}}+x^{\frac{121}{8}}+ .\,.\bigr)
\end{gathered}
\]
Nun findet man aus 6 nach dem, was zwischen 22 und 23 gezeigt ist
\[
\begin{gathered}
1+x^4y\,.\,1+x^{10}y\,.\,1+x^{16}y\,.\,.\,.\,
\\[8pt]\quad{}1+\frac{xx}{y}\,.\,1+\frac{x^8}{y}\,.\,1+\frac{x^{14}}{y}\,.\,.\,.\,\\[8pt]\quad{}
1-x^6\,.\,1-x^{12}\,.\,.\,.\\[8pt]
\qquad=\\[8pt]\quad{}1+\frac{xx}{y}+x^4y+\frac{x^{10}}{yy}+x^{14}yy+ .\,.\,.
\end{gathered}
\]
Also
\[
\begin{gathered}
1+x^4i\,.\,1-x^{10}i\,.\,1+x^{16}i\,.\,.\,.\,
\\[8pt]\quad{}1-\frac{xx}{i}\,.\,1+\frac{x^8}{i}\,.\,1-\frac{x^{14}}{i}\,.\,.\,.\,\\[8pt]\quad{}
1+x^6\,.\,1-x^{12}\,.\,1+x^{18}\,.\,.\,.\\[8pt]
\qquad=1+ixx+ix^4+x^{10}+x^{14}+ .\,.\,.
\end{gathered}
\]
Daher die Zerlegung in Factoren
\[
\begin{gathered}
47.\quad (1-2x^3+2x^{12}- .\,.)\\[8pt]\quad{}-i(1-2x+2x^4- .\,.)\\[8pt]
\qquad=1-x\,.\,1-x^3\,.\,1-x^5\,.\,.\,.\,\\[8pt]\quad{}
1-i\,.\,1+x^3\,.\,1-x^6\,.\,1+x^9\,.\,.\,.\\[8pt]
\qquad\quad\times1+ix\,.\,1+ixx\,.\,1-ix^4\,.\,1-ix^5\,.\,1+ix^7\,.\,1+ix^8\,.\,.\,.
\end{gathered}
\]
